\newpage
\section{Scope and Limitations}
\label{ch:scope-limitations}

This study focuses on the automated classification of rice pollen viability from stained microscopic images. It will use images from six rice genotypes: IR68652-3B-8-1, Sadri, CSR28, IR67075-2B-5-2, Rc222, and Pokkali, subject to permission to access the dataset used by Go et al.~\cite{go2019}. If additional images are required, samples will be collected from the same genotypes using comparable staining and imaging procedures. The study will fine-tune a pretrained YOLOv8n-seg instance-segmentation model to identify individual pollen grains and classify them as viable or nonviable. Its performance will be evaluated on images with and without lumped pollen and compared with the image-analysis protocol of Go et al.~\cite{go2019}, which uses edge detection, dilation, region filling, and $k$-means clustering.

The evaluation will focus on the model's ability to detect and classify individual grains, the effect of overlapping or clustered pollen on performance, and the agreement between automated viability estimates and manual counts. Stratified five-fold cross-validation will be used to evaluate the methods, with precision, recall, and F1-score used for detection and classification, and mean absolute error (MAE) and Pearson correlation used to assess the agreement of estimated viability percentages with manual counts. The comparison will be limited to the proposed YOLOv8n-seg model and the reimplemented protocol of Go et al.; other deep-learning architectures will not be evaluated in this study.

Several limitations may affect the findings. First, the study depends on access to the original dataset, and the number and diversity of available images may limit the range of conditions represented. The images cover six rice genotypes and may not represent all rice varieties, staining procedures, microscopes, or image-acquisition conditions. Second, the model's performance will depend on the quality and consistency of grain annotations, especially when individual grains overlap and their boundaries are difficult to identify. Third, the baseline comparison may be affected by implementation choices where details of the original protocol are unavailable; the assignment of its $k$-means clusters to viable and nonviable classes will therefore be checked against the annotated images. Finally, the study evaluates image-based, stain-based classification and estimated viability percentages. It will not directly measure pollen germination, fertilization success, crop yield, or the model's performance in field conditions. These limitations will be considered when interpreting the results and determining the extent to which the proposed approach can be generalized.
